\subsection{El clasificador identifica la derivada de Radon--Nikodym}\label{sec:clasificador}

Un clasificador calibrado que distingue el grupo objetivo del canal donante da $w$ salvo una constante, y el conjunto donde ese clasificador está seguro, $\{p=1\}$, es exactamente la parte del grupo objetivo sin soporte.

\paragraph{Construcción.} Juntamos ambas poblaciones con una etiqueta: $(\tilde X,Z)$ con $Z\sim\mathrm{Bernoulli}(\pi)$, $\pi\in(0,1)$, $\tilde X\mid Z=1\sim P$ y $\tilde X\mid Z=0\sim Q$. La ley de $\tilde X$ es la mezcla $M=\pi P+(1-\pi)Q$, y el clasificador ideal es una versión $p(x)$ de $\PP(Z=1\mid\tilde X=x)$. En la práctica $\pi$ es la fracción del grupo objetivo en la muestra combinada; no afecta al estimando.

\begin{lema}\label{lem:mezcla}
$P\ll M$ y $Q\ll M$, con $dP/dM=p/\pi$ y $dQ/dM=(1-p)/(1-\pi)$, $M$-casi seguro.
\end{lema}
\begin{proof}
Si $M(A)=0$, entonces $\pi P(A)+(1-\pi)Q(A)=0$ y, como ambos coeficientes son positivos, $P(A)=Q(A)=0$. Además, para todo $A\in\A$, $\pi P(A)=\PP(Z=1,\tilde X\in A)=\int_A p\dd M$ por definición de esperanza condicionada; luego $p/\pi$ es una versión de $dP/dM$. El caso de $Q$ es idéntico con $1-p$.
\end{proof}

\begin{teorema}[Cociente de odds]\label{teo:odds}
Tomemos una versión de $p$ con valores en $[0,1]$ y sea $r=\frac{1-\pi}{\pi}\,\frac{p}{1-p}$ sobre $\{p<1\}$ y $r=0$ sobre $\{p=1\}$. Entonces $Q(\{p=1\})=0$ y, para todo $A\in\A$, $\int_A r\dd Q=P(A\cap\{p<1\})$. En particular, bajo (A3), $r=dP/dQ$ $Q$-casi seguro.
\end{teorema}
\begin{proof}
Por el Lema~\ref{lem:mezcla}, $Q(\{p=1\})=\int_{\{p=1\}}\frac{1-p}{1-\pi}\dd M=0$. Para $A\in\A$, por la regla de la cadena,
\[
\int_A r\dd Q=\int_{A\cap\{p<1\}}\frac{1-\pi}{\pi}\,\frac{p}{1-p}\cdot\frac{1-p}{1-\pi}\dd M=\int_{A\cap\{p<1\}}\frac{p}{\pi}\dd M=P(A\cap\{p<1\}).
\]
Bajo (A3), $P(\{p=1\})=0$ porque $Q(\{p=1\})=0$, así que $\int_A r\dd Q=P(A)$ para todo $A$.
\end{proof}

Combinado con la invarianza del Teorema~\ref{teo:identificacion}, bastan los odds $p/(1-p)$: la constante $(1-\pi)/\pi$ se cancela. Con un clasificador de log-odds $F$, el peso es $e^{F}$. El teorema exige que $p$ sea la probabilidad condicional verdadera: un clasificador bien ordenado pero mal calibrado ordena bien a las unidades y aun así da pesos sesgados.

\paragraph{Ejemplo~\ref{ej:juguete} (continuación).}
Con $\pi=0,2$, el clasificador ideal vale $p=(0,048;\,0,200;\,0,429)$ en A, B y C, con odds $(0,05;\,0,25;\,0,75)=w/4$, y el estimador autonormalizado da $1,3\,\%$. El clasificador $\tilde p=p^2$ ordena igual (mismo AUC), pero sus odds $(0,0023;\,0,042;\,0,225)$ dan $0,89\,\%$, un sesgo del 32\,\%. La calibración importa, no solo la discriminación.


\begin{teorema}[Soporte y descomposición de Lebesgue]\label{teo:lebesgue}
Sin suponer (A3):
\begin{enumerate}
  \item $P\ll Q$ si y solo si $P(\{p=1\})=0$.
  \item La descomposición de Lebesgue de $P$ respecto de $Q$ es $P=\Pac+\Ps$, con $\Pac(A)=P(A\cap\{p<1\})$, $\Ps(A)=P(A\cap\{p=1\})$ y $d\Pac/dQ=r$.
  \item Sea $\varepsilon=P(\{p=1\})$. Si $\varepsilon<1$, (A2) se cumple sobre $\{p<1\}$ y $\int_{\{p<1\}}\E[|Y|\mid X=x,D=0]\,P(dx)<\infty$, la ponderación autonormalizada con $r$ recupera
  \[
  \frac{\E[r(X)\,Y\mid D=0]}{\E[r(X)\mid D=0]}=\thac:=\E\bigl[Y(0)\mid S=1,\ p(X)<1\bigr],
  \]
  y $\theta=(1-\varepsilon)\thac+\varepsilon\,\ths$, con $\ths=\E[Y(0)\mid S=1,\ p(X)=1]$, definido cuando $\varepsilon>0$, no identificado.
\end{enumerate}
\end{teorema}
\begin{proof}
(2) Por el Teorema~\ref{teo:odds}, $\Pac$ tiene densidad $r$ respecto de $Q$, luego $\Pac\ll Q$; $\Ps$ está concentrada en $\{p=1\}$, que es $Q$-nulo, luego $\Ps\perp Q$; por la unicidad de la descomposición de Lebesgue, esta es la descomposición. El conjunto $\{p=1\}$ depende de la versión de $p$ solo módulo conjuntos $M$-nulos, que son nulos para $P$ y para $Q$. (1) $P\ll Q$ si y solo si $\Ps=0$. (3) Por el Teorema~\ref{teo:odds}, $\E[r(X)\mid D=0]=1-\varepsilon$; repitiendo los pasos 3 y 4 del Teorema~\ref{teo:identificacion} con $\Pac$, el numerador vale $\int_{\{p<1\}}m_0\dd P=(1-\varepsilon)\thac$. La relación con $\theta$ es la ley de probabilidad total.
\end{proof}

\begin{corolario}[Cotas de identificación parcial]\label{cor:cotas}
Con las hipótesis del Teorema~\ref{teo:lebesgue}(3), si $Y\in[0,1]$, entonces $(1-\varepsilon)\thac\le\theta\le(1-\varepsilon)\thac+\varepsilon$, y el sesgo de la ponderación es $\thac-\theta=\varepsilon(\thac-\ths)$.
\end{corolario}
La anchura del intervalo es $\varepsilon$, que puede superar al propio $\theta$ cuando el resultado es raro: si al Ejemplo~\ref{ej:juguete} se añade un perfil sin soporte con $P=(0,1;\,0,3;\,0,5;\,0,1)$, entonces $\varepsilon=0,1$, $\thac=1,39\,\%$ y el intervalo es $[1,25\,\%;\,11,25\,\%]$. Solo es informativo con un supuesto adicional; por ejemplo, si las unidades sin soporte no responden mejor que el resto del grupo, $\ths\le\thac$, el intervalo se reduce a $[(1-\varepsilon)\thac,\thac]$, aquí $[1,25\,\%;\,1,39\,\%]$, y el estimador ponderado es una cota superior de $\theta$.

\begin{corolario}[Estimador doblemente robusto sin soporte]\label{cor:drsinsoporte}
Bajo las hipótesis del Teorema~\ref{teo:lebesgue}, sea $\bar m$ $P$-integrable con $\bar m=m_0$ $P$-casi seguro sobre $\{p<1\}$. Con los pesos $r$ (el último término es $0$ cuando $\varepsilon=0$),
\[
\frac{\E[r(X)(Y-\bar m(X))\mid D=0]}{\E[r(X)\mid D=0]}+\int\bar m\dd P=(1-\varepsilon)\thac+\varepsilon\,\E_P\bigl[\bar m(X)\mid p(X)=1\bigr].
\]
\end{corolario}
\begin{proof}
Por el Teorema~\ref{teo:odds}, el numerador del primer sumando es $\int_{\{p<1\}}(m_0-\bar m)\dd P=0$. El segundo se parte en $\{p<1\}$, donde integra $(1-\varepsilon)\thac$, y $\{p=1\}$, donde integra $\varepsilon\,\E_P[\bar m\mid p=1]$.
\end{proof}

A diferencia del estimador ponderado, que se queda con la parte con soporte, el doblemente robusto rellena la parte sin soporte con la extrapolación del modelo de resultado. Recupera $\theta$ si y solo si esa extrapolación acierta en media, $\E_P[\bar m\mid p=1]=\ths$, que es un supuesto adicional no contrastable. \citet{manamkoong2025} proponen hacerla explícita bajo un supuesto de suavidad y acotar su efecto con un análisis de sensibilidad.

\paragraph{Lectura práctica.} En muestra finita nunca se observa $\hat p=1$ exactamente. Un diagnóstico natural es la masa estimada $\hat P(\hat p>1-t)$ para $t$ pequeño; un AUC cercano a 1 es una señal de alarma, porque indica que $P$ y $Q$ están cerca de ser mutuamente singulares, no que el clasificador sea bueno. Por debajo de ese extremo, el AUC no es una medida monótona del solapamiento: cuando el grupo objetivo es un corte en una puntuación, es alto incluso con buen solapamiento (sección~\ref{sec:diseno}).

\subsection{Estimadores}\label{sec:estimadores}

\paragraph{Muestreo.} Observamos una muestra combinada i.i.d. $O_i=(Z_i,X_i,(1-Z_i)Y_i)$, $i=1,\dots,n$, con $\PP(Z=1)=\pi$, $X\mid Z=1\sim P$ y $X\mid Z=0\sim Q$; para $\theta$ solo hace falta el resultado de las unidades donantes. Cuando una unidad (un cliente) genera varias observaciones (llamadas), $O_i$ es el bloque completo de esa unidad y todo lo que sigue se aplica a nivel de unidad. Suponemos que el canal y la pertenencia al grupo objetivo son propiedades de la unidad, constantes entre sus observaciones, y que $P$, $Q$, $m_0$ y $\theta$ se definen a nivel de observación, de modo que cada unidad pesa según su número de observaciones. Si ese número depende de $Y(0)$ dado $X$, los estimandos por unidad y por observación, y las versiones correspondientes de (A2), difieren.

\paragraph{Estimadores.} Con $\hat o=\hat p/(1-\hat p)$ y $n_1$ el número de observaciones del grupo objetivo,
\[
\hat\theta_{\mathrm{H}}=\frac{\sum_{i:Z_i=0}\hat o(X_i)\,Y_i}{\sum_{i:Z_i=0}\hat o(X_i)},\qquad
\hat\theta_{\mathrm{DR}}=\frac{\sum_{i:Z_i=0}\hat o(X_i)\,(Y_i-\hat m(X_i))}{\sum_{i:Z_i=0}\hat o(X_i)}+\frac1{n_1}\sum_{i:Z_i=1}\hat m(X_i).
\]
El primero es el estimador ponderado de Hájek. El segundo es el doblemente robusto (DR), la versión muestral de $\psi$ en la Proposición~\ref{prop:dr} con el término de pesos normalizado. Ambos usan \emph{cross-fitting}: las unidades se reparten en $K$ bloques y, para las observaciones del bloque $k$, $\hat p$ y $\hat m$ se entrenan con los demás.

\paragraph{Diagnósticos y evidencia.} De un clasificador ajustado se leen habitualmente tres diagnósticos: la masa sin soporte estimada $\hat\varepsilon=\hat P(\hat p>1-t)$ para un $t$ pequeño (Teorema~\ref{teo:lebesgue}); el AUC; y el tamaño muestral efectivo de Kish \citeyearpar{kish1965}, $\neff=(\sum_i\hat w_i)^2/\sum_i\hat w_i^2$ sobre las $n_0$ observaciones donantes, cuyo cociente $\neff/n_0$ estima $1/(1+\chi^2(P\|Q))$ (sección~\ref{sec:inferencia}). Las secciones~\ref{sec:canal} y~\ref{sec:instrumentos} estudian qué muestran estos diagnósticos cuando los canales registran las covariables de forma distinta. En todo el texto, $z$ denota (media $-$ predicción)/(error de Monte Carlo) entre réplicas, y una predicción está \emph{registrada} si se subió, con fecha, al repositorio público (sección~\ref{sec:diseno}) antes de la ejecución correspondiente.
